% ============================================================
%  APUNTES — Signals en Angular (Survey Maker)
%  Compilar con: pdflatex apuntes-signals.tex
% ============================================================
\documentclass[12pt, a4paper]{article}

% --- Codificación y fuentes (pdfLaTeX) ----------------------
\usepackage[utf8]{inputenc}
\usepackage[T1]{fontenc}
\usepackage{lmodern}

% --- Idioma -------------------------------------------------
\usepackage[spanish]{babel}

% --- Márgenes -----------------------------------------------
\usepackage[top=2.5cm, bottom=2.5cm, left=2.8cm, right=2.8cm]{geometry}

% --- Colores ------------------------------------------------
\usepackage[dvipsnames, table]{xcolor}
\definecolor{azulSignal}{RGB}{37, 99, 235}
\definecolor{verdeOk}{RGB}{22, 163, 74}
\definecolor{rojoError}{RGB}{220, 38, 38}
\definecolor{grisClaro}{RGB}{243, 244, 246}
\definecolor{grisCodigo}{RGB}{31, 41, 55}
\definecolor{amarillo}{RGB}{251, 191, 36}
\definecolor{morado}{RGB}{124, 58, 237}

% --- Código fuente ------------------------------------------
\usepackage{listings}
\lstdefinestyle{ts}{
  language        = [Sharp]C,           % lo más cercano a TypeScript
  basicstyle      = \ttfamily\small,
  keywordstyle    = \color{azulSignal}\bfseries,
  commentstyle    = \color{verdeOk}\itshape,
  stringstyle     = \color{rojoError},
  backgroundcolor = \color{grisClaro},
  frame           = single,
  framerule       = 0.5pt,
  rulecolor       = \color{gray!40},
  numbers         = left,
  numberstyle     = \tiny\color{gray},
  numbersep       = 8pt,
  tabsize         = 2,
  breaklines      = true,
  breakatwhitespace = true,
  showstringspaces = false,
  xleftmargin     = 1.5em,
  morekeywords    = {signal, computed, effect, inject, readonly,
                     const, let, async, of, map, return, true, false,
                     null, undefined, string, number, boolean, void},
}
\lstset{style=ts}

% --- Cajas decorativas --------------------------------------
\usepackage[most]{tcolorbox}

\tcbset{
  analogia/.style={
    colback   = blue!5,
    colframe  = azulSignal,
    fonttitle = \bfseries,
    title     = Analogía,
    arc       = 4pt,
  },
  importante/.style={
    colback   = yellow!10,
    colframe  = amarillo,
    fonttitle = \bfseries,
    arc       = 4pt,
  },
  regla/.style={
    colback   = green!5,
    colframe  = verdeOk,
    fonttitle = \bfseries,
    arc       = 4pt,
  },
}

% --- Hiperlinks ---------------------------------------------
\usepackage{hyperref}
\hypersetup{
  colorlinks = true,
  linkcolor  = azulSignal,
  urlcolor   = morado,
}

% --- Tablas -------------------------------------------------
\usepackage{booktabs}
\usepackage{array}

% --- Encabezado y pie de página ----------------------------
\usepackage{fancyhdr}
\pagestyle{fancy}
\fancyhf{}
\rhead{\textcolor{gray}{\small Signals — Angular Moderno}}
\lhead{\textcolor{gray}{\small Survey Maker}}
\cfoot{\textcolor{gray}{\thepage}}

% --- Gráficos (diagramas ASCII con tikz) --------------------
\usepackage{tikz}
\usetikzlibrary{arrows.meta, shapes, positioning}

% ============================================================
\begin{document}

% --- Portada ------------------------------------------------
\begin{titlepage}
  \centering
  \vspace*{3cm}
  {\Huge\bfseries\color{azulSignal} Apuntes de Angular}\\[0.5cm]
  {\Large\color{grisCodigo} Signals — El Sistema Reactivo}\\[1cm]
  \rule{10cm}{1pt}\\[1cm]
  {\large Proyecto: \textbf{Survey Maker}}\\[0.3cm]
  {\large\texttt{D:/Codigo Abierto/GenEncuesta/repo/survey-maker}}\\[1cm]
  {\normalsize Aldo Zetina — UNAM Ciencias de la Computación}\\[0.3cm]
  {\normalsize\color{gray} Mayo 2026}\\[2cm]
  \begin{tcolorbox}[importante, title=Cómo compilar este archivo]
    Compatible con \textbf{pdfLaTeX} y \textbf{XeLaTeX}.\\
    Comando: \texttt{pdflatex apuntes-signals.tex}\\
    En Overleaf: compilador por defecto funciona.
  \end{tcolorbox}
\end{titlepage}

\tableofcontents
\newpage

% ============================================================
\section{La Idea Central — Sin Código}

Antes de ver código, entiende el concepto con una historia.

\begin{tcolorbox}[analogia, title=La Pizarra Mágica]
Imagina que tienes una \textbf{pizarra mágica} en tu cuarto. Escribes un número en ella, digamos \textit{``0 respuestas contestadas''}.

Afuera hay tres amigos mirando por la ventana:
\begin{itemize}
  \item \textbf{Barra de progreso} — actualiza su dibujo cuando cambia el número.
  \item \textbf{Botón de Enviar} — se activa cuando el número llega al máximo.
  \item \textbf{localStorage} — guarda una foto del número automáticamente.
\end{itemize}

Cada vez que \textbf{borras el número y escribes uno nuevo}, los tres amigos lo ven al instante. Tú no tienes que avisarles — ellos solos reaccionan.

\textbf{Eso es un Signal.} La pizarra mágica.
\end{tcolorbox}

% ============================================================
\section{Tipo 1 — \texttt{signal()} — La Pizarra Mágica}

\subsection{Qué es}

Un \texttt{signal()} crea una variable reactiva. Cuando su valor cambia,
Angular automáticamente actualiza todo lo que depende de él.

\subsection{En el proyecto}

Archivo: \texttt{src/app/core/services/survey-state.service.ts}, línea 28.

\begin{lstlisting}
// signal() crea la "pizarra mágica".
// Todo empieza vacío {} porque nadie ha contestado nada todavía.
// Record<string, FieldAnswer> = un diccionario donde:
//   - la LLAVE es el ID del campo  (ej: "nombre", "calificacion")
//   - el VALOR es la respuesta     (ej: { value: "Juan", isValid: true })
readonly answers = signal<Record<string, FieldAnswer>>({});
//                                                      ^
//                                                      valor inicial — pizarra vacía
\end{lstlisting}

\subsection{Cómo leer y escribir un Signal}

\begin{lstlisting}
// LEER — siempre con paréntesis ()
const todasLasRespuestas = this.answers();
//                                      ^
//                          los () = "dame lo que hay en la pizarra"

// ESCRIBIR — con .set() — borra todo y pone el nuevo valor
this.answers.set({ nombre: { value: "Juan", isValid: true } });

// ACTUALIZAR — con .update() — modifica sin borrar el resto
this.answers.update(prev => ({
//                    ^
//                    prev = lo que había antes en la pizarra
  ...prev,           // copia todo lo anterior (no lo borres)
  [fieldId]: {       // solo cambia ESTE campo específico
    value: "nuevo valor",
    isValid: true,
    isDirty: true,
  }
}));
\end{lstlisting}

\begin{tcolorbox}[regla, title=Regla de Oro — Signals]
  Siempre leer con \texttt{signal()}, nunca sin paréntesis.\\
  \texttt{this.answers} → objeto Signal (no el valor).\\
  \texttt{this.answers()} → el valor real (el diccionario).
\end{tcolorbox}

% ============================================================
\section{Tipo 2 — \texttt{computed()} — El Amigo Matemático}

\subsection{Qué es}

Un \texttt{computed()} es una variable que se \textbf{calcula sola} a partir de
otros Signals. Cuando cualquier Signal que lee cambia, el \texttt{computed}
recalcula automáticamente.

\begin{tcolorbox}[analogia, title=Celda de Excel con Fórmula]
  Imagina \texttt{C1 = A1 + B1} en Excel.\\
  Cuando cambias \texttt{A1} o \texttt{B1}, \texttt{C1} se actualiza solo.\\
  Tú no tienes que recalcularlo a mano.\\[0.3cm]
  \texttt{progress} es \texttt{C1}. \texttt{answers} es \texttt{A1}.
\end{tcolorbox}

\subsection{Progreso — en el proyecto}

Archivo: \texttt{src/app/core/services/survey-state.service.ts}, línea 43.

\begin{lstlisting}
readonly progress = computed(() => {

  const s = this.schema();
  // "schema" es otra pizarra con la definición de la encuesta.
  // Si no hay encuesta cargada todavía → regresa 0%
  if (!s) return 0;

  // Aplanamos todas las secciones en un solo array de campos.
  // Imagina: [campo1, campo2, campo3, campo4, campo5]
  const allFields = s.sections.flatMap(sec => sec.fields);

  // Solo nos importan los campos OBLIGATORIOS (los que tienen *)
  const requiredFields = allFields.filter(f => f.validation?.required);

  // Si no hay campos obligatorios → ya terminaste, 100%
  if (requiredFields.length === 0) return 100;

  // Contamos cuántos campos obligatorios ya contestó el usuario.
  // Un campo cuenta como "contestado" SOLO si:
  //   isValid = true  (el valor es correcto)
  //   isDirty = true  (el usuario realmente lo tocó)
  const answeredRequired = requiredFields.filter(f => {
    const ans = this.answers()[f.id];
    //               ^
    //               Lee "answers" Signal — aquí crea la dependencia.
    //               Si "answers" cambia → este computed recalcula solo.
    return ans?.isValid && ans.isDirty;
  });

  // Fórmula: (contestados / total) × 100, redondeado a entero.
  // Ejemplo: 3 de 5 campos = (3/5) × 100 = 60%
  return Math.round((answeredRequired.length / requiredFields.length) * 100);
});
\end{lstlisting}

\subsection{Validez — en el proyecto}

Archivo: \texttt{src/app/core/services/survey-state.service.ts}, línea 72.

\begin{lstlisting}
// Este computed responde UNA pregunta:
// ¿puede el usuario enviar la encuesta?
readonly isValid = computed(() =>

  // .every() = "¿TODOS los campos cumplen la condición?"
  // Si UN SOLO campo falla → regresa false → botón bloqueado
  this.schema()?.sections
    .flatMap(s => s.fields)
    .every(f => {

      // Si el campo NO es obligatorio → lo ignoramos, cuenta como válido
      if (!f.validation?.required) return true;

      // Si ES obligatorio → revisamos si tiene respuesta válida
      const ans = this.answers()[f.id];
      //                ^
      //                Otra lectura de "answers" → dependencia mágica
      return ans?.isValid ?? false;
      //                   ^
      //                   ?? false = "si no existe respuesta → false"
    }) ?? false
    //   ^
    //   ?? false = "si no hay schema todavía → false"
);
\end{lstlisting}

% ============================================================
\section{Tipo 3 — \texttt{effect()} — El Guardián Automático}

\subsection{Qué es}

Un \texttt{effect()} ejecuta código cada vez que un Signal cambia.
No retorna nada — solo realiza \textbf{efectos secundarios}: guardar datos,
hacer log, llamar APIs externas.

\begin{tcolorbox}[analogia, title=Alarma Automática]
  Es como configurar una alarma que suena \textit{cada vez} que alguien
  entra a tu cuarto.\\
  La configuras una vez y olvidas — ella solita se activa cuando
  ocurre el evento.
\end{tcolorbox}

\subsection{En el proyecto}

Archivo: \texttt{src/app/core/services/survey-state.service.ts}, línea 88.

\begin{lstlisting}
effect(() => {

  // untracked() = "lee este Signal pero NO lo vigiles"
  // Sin untracked, el effect correría cada vez que schema cambia también.
  // Solo queremos reaccionar cuando ANSWERS cambia.
  const s = untracked(() => this.schema());

  // ESTA línea SÍ crea la dependencia.
  // Cuando "answers" cambia → este effect corre automáticamente.
  const answers = this.answers();

  // Si no hay encuesta cargada → no guardes nada
  if (!s) return;

  try {
    // Autoguardado en localStorage — como Google Docs guardando tu documento.
    // Clave única por encuesta para no mezclar datos.
    // Ej: "survey-state-ejemplo-feedback" → { campo1: {...} }
    localStorage.setItem(
      `survey-state-${s.id}`,   // llave única por encuesta
      JSON.stringify(answers)    // convertimos objeto a texto para guardar
    );
  } catch {
    // localStorage puede fallar en modo privado o si está lleno.
    // No hacemos nada — la encuesta funciona igual, solo sin autoguardado.
  }

}, { injector: this.injector });
\end{lstlisting}

% ============================================================
\section{Tipo 4 — \texttt{computed()} para Lógica Condicional}

\subsection{Qué hace}

Decide si un campo de la encuesta se muestra o no, basándose en
la respuesta de otro campo.

Archivo: \texttt{src/app/features/survey/field-renderer.component.ts}, línea 156.

\begin{lstlisting}
// Decide si este campo es visible en pantalla.
// Ejemplo real: "el campo de comentarios solo aparece
//               si elegiste 'Malo' en la calificación"
readonly isVisible = computed(() => {
  const condition = this.field().condition;

  // Si el campo no tiene condición → siempre visible
  if (!condition) return true;

  // Buscamos la respuesta del campo "controlador"
  // condition.dependsOn = ID del campo que controla a este
  const answer = this.answers()[condition.dependsOn];
  const expected = condition.equals;

  // condition.equals puede ser un array para varios valores válidos.
  // Ejemplo: equals: ["malo", "regular"]
  // → el campo aparece si eligieron "malo" O "regular"
  if (Array.isArray(expected)) {
    return expected.includes(answer as string);
  }

  // Comparación simple para un solo valor esperado
  return answer === expected;
});
\end{lstlisting}

\begin{tcolorbox}[importante, title=Importante — \texttt{@if} en el Template]
Cuando \texttt{isVisible()} regresa \texttt{false}, Angular \textbf{elimina el componente
del DOM completamente}. No es \texttt{display: none} — el componente
no existe en memoria y no ejecuta código.
\end{tcolorbox}

% ============================================================
\section{Los Tres Tipos — Resumen Visual}

\begin{center}
\begin{tabular}{>{\bfseries}p{3cm} p{4.5cm} p{5cm}}
\toprule
Tipo & Para qué & Ejemplo en el proyecto \\
\midrule
\texttt{signal()} & Guardar estado que cambia & \texttt{answers}, \texttt{schema}, \texttt{isSubmitted} \\
\addlinespace
\texttt{computed()} & Calcular algo a partir de Signals & \texttt{progress}, \texttt{isValid}, \texttt{isVisible} \\
\addlinespace
\texttt{effect()} & Reaccionar con efectos secundarios & Guardar en \texttt{localStorage} \\
\bottomrule
\end{tabular}
\end{center}

% ============================================================
\section{Flujo Completo — Lo Que Pasa Con Cada Tecla}

Cuando el usuario escribe una letra en cualquier campo:

\begin{enumerate}
  \item El browser dispara el evento \texttt{(input)}.
  \item \texttt{TextFieldComponent} emite \texttt{valueChange("texto")}.
  \item \texttt{SurveyComponent} llama \texttt{state.setAnswer('campo-id', 'texto')}.
  \item \texttt{SurveyStateService} valida el valor y llama \texttt{answers.update(...)}.
  \item \textbf{Automáticamente:}
    \begin{itemize}
      \item \texttt{progress} computed recalcula → la barra de progreso sube.
      \item \texttt{isValid} computed recalcula → el botón se habilita si corresponde.
      \item El \texttt{effect()} detecta el cambio → guarda en \texttt{localStorage}.
      \item \texttt{isVisible} computed de cada campo reevalúa → campos condicionales aparecen/desaparecen.
    \end{itemize}
\end{enumerate}

\begin{tcolorbox}[regla, title=El Principio Más Importante]
  \textbf{Tú nunca avisas manualmente.}\\
  Solo cambias el Signal con \texttt{.set()} o \texttt{.update()}.\\
  Todo lo demás — la UI, el progreso, el localStorage — se actualiza solo.
\end{tcolorbox}

% ============================================================
\section{Inmutabilidad — Por Qué Nunca Mutamos Directo}

En todo el proyecto verás el patrón \texttt{...prev}. Aquí está el porqué:

\begin{lstlisting}
// INCORRECTO — misma referencia, Angular NO detecta el cambio
prev[fieldId] = nuevoValor;
this.answers.set(prev);    // Angular ve el mismo objeto → no re-renderiza

// CORRECTO — nueva referencia, Angular SÍ detecta el cambio
this.answers.update(prev => ({
  ...prev,                  // copia todos los campos anteriores
  [fieldId]: nuevoValor     // sobreescribe solo el que cambió
}));
// Resultado: objeto NUEVO en memoria → Angular detecta → UI se actualiza
\end{lstlisting}

\begin{tcolorbox}[analogia, title=Por Qué Importa la Referencia]
  Angular compara referencias de objetos, no su contenido.\\
  Es como comparar la \textit{dirección de una casa} en lugar del \textit{contenido}.\\
  Si la dirección no cambia → Angular asume que nada cambió → no actualiza.\\
  Con \texttt{\{...prev, campo: nuevo\}} → dirección nueva → Angular actualiza.
\end{tcolorbox}

% ============================================================
\section{La Analogía Final — La Pizzería}

\begin{center}
\begin{tabular}{>{\bfseries}p{4cm} p{8cm}}
\toprule
Concepto Angular & Analogía \\
\midrule
\texttt{signal(answers)} & La orden escrita en papel \\
\texttt{computed(progress)} & La pantalla que dice ``tu pizza: 60\% lista'' \\
\texttt{computed(isValid)} & El semáforo que dice ``listo para hornear'' \\
\texttt{effect()} & El sistema que imprime el recibo automáticamente \\
\texttt{.update()} & El cocinero modifica la orden \\
Las pantallas & Los \texttt{computed()} que reaccionan solos \\
\bottomrule
\end{tabular}
\end{center}

Cuando el cocinero modifica la orden, \textbf{todas las pantallas cambian solas}.
No tiene que ir pantalla por pantalla avisando.

% ============================================================
\section{Archivos Clave para Repasar}

\begin{center}
\begin{tabular}{p{7cm} p{5cm}}
\toprule
\textbf{Archivo} & \textbf{Qué estudiar} \\
\midrule
\texttt{core/services/survey-state.service.ts} & \texttt{signal}, \texttt{computed}, \texttt{effect} \\
\texttt{features/survey/field-renderer.component.ts} & \texttt{isVisible} computed \\
\texttt{features/fields/text-field.component.ts} & \texttt{input()}, \texttt{output()} \\
\texttt{features/survey/survey.component.ts} & cómo conecta todo \\
\bottomrule
\end{tabular}
\end{center}

% ============================================================
\vfill
\begin{center}
  \textcolor{gray}{\small Apuntes generados en sesión de estudio — Mayo 2026}\\
  \textcolor{gray}{\small Proyecto Survey Maker — Angular 19 Signal-First}
\end{center}

\end{document}
